\begin{corollary}[6.7.8]
\label{IV.2.s6.7.8-fr}
Sea $k$ un cuerpo; para un $k$-preesquema localmente noetheriano $X$ y un
punto $x\in X$, se designa por $\boldsymbol{P}(X,x)$ una de las propiedades siguientes:
\begin{enumerate}
\item[\textup{(i)}] $\operatorname{coprof}(\mathcal O_x)\leq n$.
\item[\textup{(ii)}] $\mathcal O_x$ es un anillo de Cohen-Macaulay.
\item[\textup{(iii)}] $X$ verifica la propiedad $(S_n)$ en el punto $x$.
\item[\textup{(iv)}] $X$ es geométricamente regular en el punto $x$.
\item[\textup{(v)}] $X$ posee la propiedad $(R_n)$ geométrica en el punto $x$.
\item[\textup{(vi)}] $X$ es geométricamente normal en el punto $x$.
\item[\textup{(vii)}] $X$ es geométricamente reducido (\textup{es decir,} separable) en el punto $x$.
\end{enumerate}
Sea $k'$ una extensión de $k$; se supone, o bien que $k'$ es de tipo finito,
o bien que $X$ es localmente de tipo finito sobre $k$, de modo que
$X'=X\otimes_k k'$ es localmente noetheriano. Sea $x'$ un punto de $X'$ del que
$x$ es la proyección en $X$. Entonces, para que $\boldsymbol{P}(X,x)$ sea verdadera, es
necesario y suficiente que $\boldsymbol{P}(X',x')$ lo sea.
\end{corollary}

\begin{proof}
Esto ya se ha visto para la propiedad (vii) (4.6.11), y para las propiedades
(i), (ii) y (iii) (6.7.1). Para (iv), (v) y (vi), resulta de (6.7.7):
en efecto, el hecho de que la condición sea necesaria resulta de la equivalencia de los
criterios \emph{a)} y \emph{c)}, y de la equivalencia de \emph{c)} y \emph{d)} cuando $X$ es localmente
de tipo finito sobre $k$. Para ver que la condición es suficiente, sea $k''$ una
extensión finita radicial de $k$; siempre se pueden considerar $k'$ y $k''$
como subextensiones de una extensión $K$ de $k$; pongamos
$X''=X\otimes_k k''$, $X_0=X\otimes_k K$, y observemos que, puesto que $k''$ es una
extensión radicial de $k$, hay \emph{un solo} punto $x''$ de $X''$ situado sobre
$x$ (\textbf I, 3.5.7 y 3.5.8). Sea entonces $x_0$ un punto cualquiera de $X_0$
situado sobre $x'$; si $\boldsymbol{P}(X',x')$ es verdadera, ocurre lo mismo con $\boldsymbol{P}(X_0,x_0)$
en virtud de (6.7.7, \emph{c)} y \emph{d)}) (pues, si $k'$ es una extensión de $k$ de tipo
finito, se puede suponer que ocurre lo mismo con $K$, y si $X$ es localmente
de tipo finito sobre $k$, $X'$ es localmente de tipo finito sobre $k'$). Se
deduce entonces de (6.7.4) que la propiedad $\boldsymbol{Q}(k'')$ es verdadera (con las notaciones de
(6.7.7)), y por tanto $\boldsymbol{P}(X,x)$ es verdadera por (6.7.7, \emph{b)}).
\end{proof}
